% ============================================================================
% Section 6.7 — Transformations on the Coordinate Plane
% State-specific: Indiana, Virginia
% ============================================================================
\section{Transformations on the Coordinate Plane}

\begin{stepGoal}
\begin{itemize}
    \item Identify and perform translations, reflections, and rotations
    \item Describe how each transformation changes coordinates
\end{itemize}
\end{stepGoal}

\begin{stepVocabBox}
    \vocabItem{Translation}{A slide --- move every point the same distance and direction.}
    \vocabItem{Reflection}{A flip --- mirror a figure across a line (axis).}
    \vocabItem{Rotation}{A turn --- spin a figure around a point.}
\end{stepVocabBox}

\begin{stepsCard}[How to Transform a Point on the Coordinate Plane]
    \stepItem{Identify the \textbf{type} of transformation: translation, reflection, or rotation.}
    \stepItem{Apply the \textbf{coordinate rule}:\par
    \begin{tabular}{@{}ll@{}}
    Translation right $a$, up $b$: & $(x + a,\; y + b)$ \\
    Reflect over $x$-axis: & $(x,\; -y)$ \\
    Reflect over $y$-axis: & $(-x,\; y)$ \\
    $90°$ clockwise about origin: & $(y,\; -x)$ \\
    \end{tabular}}
    \stepItem{Write the \textbf{new coordinates} and verify the figure keeps the same size and shape.}
\end{stepsCard}

% --- Worked Example 1 ---
\begin{stepExample}{Translate $A(2, 5)$ right $3$ and down $4$.}
    \stepShow{1} This is a \textbf{translation} (slide).

    \stepShow{2} Rule: $(x + 3,\; y - 4)$. Apply: $(2 + 3,\; 5 - 4) = (5,\; 1)$.

    \stepShow{3} The new point is $A'(5, 1)$.
    \stepResult{$A'(5,\; 1)$}
\end{stepExample}

% --- Worked Example 2 ---
\begin{stepExample}{Reflect $C(4, -2)$ over the $x$-axis.}
    \stepShow{1} This is a \textbf{reflection} over the $x$-axis.

    \stepShow{2} Rule: $(x,\; -y)$. Apply: $(4,\; -(-2)) = (4,\; 2)$.

    \stepShow{3} The new point is $C'(4, 2)$.
    \stepResult{$C'(4,\; 2)$}
\end{stepExample}

\mascotStepTip{Reflecting over the $x$-axis flips the $y$-sign. Reflecting over the $y$-axis flips the $x$-sign.}

% ============================================================================
% PRACTICE
% ============================================================================
\newpage
\begin{stepPractice}{Transformations Practice}
    \resetProblems

    \stepPracticeHeader[step1Color]{Name the Transformation}
    \begin{multicols}{2}
    \prob Sliding every point $5$ units right is a \answerBlank[2.5cm].
    \answer{translation}
    \prob Flipping over the $y$-axis is a \answerBlank[2.5cm].
    \answer{reflection}
    \end{multicols}

    \stepPracticeHeader[step2Color]{Apply the Rule}
    \prob Translate $B(1, -3)$ left $4$ and up $6$. Find $B'$. \answerBlank[2.5cm]
    \answer{$B'(-3,\; 3)$}

    \prob Reflect $D(-5, 3)$ over the $y$-axis. Find $D'$. \answerBlank[2.5cm]
    \answer{$D'(5,\; 3)$}

    \prob Rotate $E(3, 1)$ by $90°$ clockwise about the origin. Find $E'$. \answerBlank[2.5cm]
    \answer{$E'(1,\; -3)$}

    \wordProblem{Triangle $PQR$ has vertices $P(1,2)$, $Q(4,2)$, $R(4,5)$. Translate it right $2$, down $3$. List the new vertices.}{~}
    \answer{$P'(3,-1)$, $Q'(6,-1)$, $R'(6,2)$}
\end{stepPractice}
